\section{Summary and roadmap} \label{sec:summary}

\pkg{econirl} unifies twelve estimators for dynamic discrete choice and inverse reinforcement learning under a single result-object contract with a shared post-estimation pipeline. The shared pipeline provides four standard error methods, identification diagnostics, profile likelihood, hypothesis testing, counterfactual simulation, and visualization. Every improvement to the inference layer benefits every estimator at once. Every new estimator added to the library inherits the inference layer for free. The user picks the estimator that fits the problem and uses the same workflow regardless.

The three illustrations in Section~\ref{sec:illustrations} demonstrate the value of the unified workflow. The same Rust bus engine replacement panel admits both a structural likelihood fit through the nested fixed point estimator and an inverse reinforcement learning fit through the maximum causal entropy estimator. The two estimators recover the same parameter vector to four decimal places, confirming the theoretical equivalence under linear features in a setting where running both estimators required only a one-line code change. The Keane-Wolpin occupational choice panel demonstrates that the neural value approximation family scales to state spaces too large for tabular methods while preserving the unified inference workflow. The cross-estimator benchmark in Table~\ref{tab:bench} demonstrates that all twelve estimators interoperate on the same data with the same outputs available for downstream analysis.

Three extensions are on the roadmap. Continuous-action problems will be supported through a Gauss-Hermite quadrature integration of the action space, which fits the existing protocol with minor changes to the choice probability computation. Dynamic games will be supported through a multi-agent extension that adds a best-response iteration to the existing inner loop. Bayesian inference will be supported through a Hamiltonian Monte Carlo sampler over the parameters, taking advantage of the JAX-computed gradients that the package already produces for maximum likelihood estimation. Each of these extensions preserves the result-object contract that defines the library.

\pkg{econirl} is open source under the MIT license. Contributions, bug reports, and feature requests are welcome at \url{https://github.com/rawatpranjal/econirl}. The package documentation is hosted at \url{https://econirl.readthedocs.io}.
